\documentclass[10pt]{article}
\usepackage[a4paper,margin=25mm]{geometry}
\usepackage{amsmath,amssymb,amsthm}
\usepackage[hidelinks]{hyperref}
\newtheorem{theorem}{Theorem}
\begin{document}
\title{Completeness of the invariant-subspace lattice: Conjecture 6825}
\author{Chengzhe Gao}\date{}\maketitle

Let $X$ be a real or complex topological vector space and let $A:X\to X$
be a linear operator. Write $\operatorname{Lat}A$ for the closed linear
subspaces $M$ with $A(M)\subseteq M$. The standard order is inclusion.
No boundedness or continuity assumption on $A$ is needed for the following
conclusion. If the algebraic convention omits topological closedness,
the same proof applies with every subset declared closed.

\begin{theorem}
The partially ordered set $(\operatorname{Lat}A,\subseteq)$ is a complete
lattice. For every family $\mathcal F\subseteq\operatorname{Lat}A$,
\[
 \bigwedge\mathcal F=\bigcap_{M\in\mathcal F}M,
 \qquad
 \bigvee\mathcal F=
 \bigcap\{U\in\operatorname{Lat}A:
                  M\subseteq U\text{ for every }M\in\mathcal F\}.
\]
The intersection over an empty family is $X$.
\end{theorem}

\begin{proof}
An arbitrary intersection of closed sets is closed. An intersection of
linear subspaces contains zero and is closed under addition and scalar
multiplication. If $x$ belongs to every member of an invariant family,
then so does $Ax$. Thus the first displayed intersection is an element
of $\operatorname{Lat}A$. It is contained in every member of $\mathcal F$,
and every common lower bound is contained in the intersection; hence it
is the greatest lower bound.

The family of upper bounds is nonempty, since it contains $X$. Its
intersection is again a closed invariant linear subspace. Every member
of $\mathcal F$ is contained in this intersection, because it is contained
in each upper bound. Conversely the intersection is contained in each
upper bound. It is therefore the least upper bound. This argument includes
the empty family, so arbitrary meets and joins exist. Inclusion is
reflexive and transitive, and antisymmetry follows from equality of the
underlying subsets. These are the complete-lattice axioms in the standard
order.
\end{proof}

\paragraph{Meaning of the source.}
The source's English and Chinese statements concern completeness in the
standard order. A complete lattice need not be totally ordered. The result
above proves existence of arbitrary infima and suprema; it makes no claim
that arbitrary invariant subspaces are comparable.

\paragraph{Formal verification.}
The self-contained Lean file represents subsets by predicates $X\to
\mathrm{Prop}$ and specifies the usual closed-set axioms of a topology.
An invariant subspace carries proofs of closedness, zero membership,
closure under the actual addition and scalar multiplication operations,
and invariance under the actual map $A$. The proof needs no vector-space
identities and is therefore proved in the stronger generality of arbitrary
operations of the same arities. Every real or complex vector space supplies
these operations; no numerical or finite-dimensional replacement is made.

Families are arbitrary predicates on the type of invariant subspaces,
including infinite and empty families. The closedness proof explicitly
converts such a family into a family of its underlying subsets and proves
equality of the two intersection predicates. Thus the closed-set axiom is
applied to the actual intersection. The carrier of the join is defined as
the intersection of all upper bounds. Extensional equality of subspaces,
all order axioms, and both universal properties are proved.

The theorem \texttt{conjecture6825} bundles these complete-lattice
properties. A discrete-topology instance covers the convention of
algebraic subspaces. A concrete instance with integers, their usual
operations, and $A(n)=2n$ additionally witnesses consistency of the more
general interface. It is not a restriction of the theorem to integers.
The project uses Lean 4.19.0 and \texttt{Std} only, with no unproved
axioms or external decision oracle; printed dependencies are the standard
logical axioms \texttt{propext} and \texttt{Quot.sound}.
\end{document}
